% BHU PYQ -- Transcribed & Corrected
% Paper: AYUVA-21: Ayurveda -- Traditional Health Care & Management Principles
% Exam: Under Graduate Semester II Examination, 2024-25 (NEP)
% Category: Value Added Course (VAC)

\documentclass[11pt,a4paper]{article}
\usepackage[a4paper,top=2.5cm,bottom=2.5cm,left=2.8cm,right=2.8cm,headheight=14pt]{geometry}
\usepackage{amsmath,amssymb}
\usepackage[T1]{fontenc}
\usepackage[utf8]{inputenc}
\usepackage{lmodern,microtype,fancyhdr,enumitem,hyperref,lastpage,parskip}

\hypersetup{
    colorlinks=true,
    urlcolor=black,
    linkcolor=black,
    pdftitle={AYUVA21: Ayurveda -- BHU Sem II 2024-25 (NEP VAC)}
}

\pagestyle{fancy}
\fancyhf{}
\renewcommand{\headrulewidth}{0.4pt}
\renewcommand{\footrulewidth}{0.4pt}
\fancyhead[L]{\small   \textit{Banaras Hindu University}}
\fancyhead[R]{\small   \textit{AYUVA-21: Ayurveda (VAC)}}
\fancyfoot[C]{\small Page  \thepage\ of \pageref{LastPage}}


\newlist{parts}{enumerate}{2}
\setlist[parts,1]{label=(\Alph*),leftmargin=*,itemsep=3pt,topsep=2pt}

\newcommand{\pts}[1]{\hfill   \textit{[#1]}}


\begin{document}


\begin{center}
  {\large\bfseries BANARAS HINDU UNIVERSITY}\\[2pt]
  {
\normalsize Under Graduate --- Semester II Examination, 2024-25 (NEP)}\\[6pt]
  \rule{0.8\linewidth}{0.5pt}\\[6pt]
  {\large\bfseries Value Added Course (VAC)}\\[4pt]
  {\large\bfseries Paper Code: AYUVA-21 : Ayurveda --- Traditional Health Care \& Management Principles}\\[4pt]
  {
\normalsize Time: 1.5 Hours\hfill Maximum Marks: 70}
\end{center}

\rule{\linewidth}{0.4pt}

\smallskip



\noindent   \textit{Instructions: Answer all 70 Multiple Choice Questions. Each question carries equal marks (1 mark each).}

\bigskip


\noindent   \textbf{Question 1.} Which factors influence health in Ayurveda?\pts{1}

\begin{parts}
  \item Diet
  \item Sleep
  \item Lifestyle
  \item All of the above
\end{parts}

\medskip


\noindent   \textbf{Question 2.} Which is not one of the four objectives of Purushartha Chatushtaya?\pts{1}

\begin{parts}
  \item Moksha
  \item Kama
  \item Dharma
  \item Vairagya
\end{parts}

\medskip


\noindent   \textbf{Question 3.} Who is the original author of Charaka-Samhita?\pts{1}

\begin{parts}
  \item Charaka
  \item Punarvasu Atreya
  \item Agnivesha
  \item Dridhabala
\end{parts}

\medskip


\noindent   \textbf{Question 4.} How many total chapters are present in Charaka-Samhita?\pts{1}

\begin{parts}
  \item 108
  \item 120
  \item 150
  \item 180
\end{parts}

\medskip


\noindent   \textbf{Question 5.} What are the primary components of "Ayu" in Ayurveda?\pts{1}

\begin{parts}
  \item Body and Soul
  \item Sensory Organs and Mind
  \item Body, Sensory Organs, Mind and Soul
  \item Soul and Consciousness
\end{parts}

\medskip


\noindent   \textbf{Question 6.} Which branch of Ayurveda deals with pediatrics?\pts{1}

\begin{parts}
  \item Kaya Chikitsa
  \item Kaumarbhritya
  \item Graha Chikitsa
  \item Shalya Tantra
\end{parts}

\medskip


\noindent   \textbf{Question 7.} What are the Panchmahabhutas?\pts{1}

\begin{parts}
  \item Five basic elements of the universe
  \item Five types of diseases
  \item Five sensory organs
  \item Five forms of meditation
\end{parts}

\medskip


\noindent   \textbf{Question 8.} The quality "Ushna" (hot) is associated with which dosha?\pts{1}

\begin{parts}
  \item Vata
  \item Pitta
  \item Kapha
  \item None of the above
\end{parts}

\medskip


\noindent   \textbf{Question 9.} Pratyaksha is a valid means of knowledge referring to:\pts{1}

\begin{parts}
  \item Inference
  \item Comparison
  \item Perception
  \item Testimony
\end{parts}

\medskip


\noindent   \textbf{Question 10.} What is Ayurveda primarily referred to as in ancient texts?\pts{1}

\begin{parts}
  \item Science of Surgery
  \item Science of Life
  \item Science of Herbs
  \item Science of Meditation
\end{parts}

\medskip


\noindent   \textbf{Question 11.} Which principle promotes moderation in physical exertion, work, rest and leisure?\pts{1}

\begin{parts}
  \item Ritucharya
  \item Sadvritta
  \item Trayopastambha
  \item Dinacharya
\end{parts}

\medskip


\noindent   \textbf{Question 12.} The three Sharirik doshas in Ayurveda are:\pts{1}

\begin{parts}
  \item Sattva, Rajas, Tamas
  \item Vata, Pitta, Kapha
  \item Agni, Soma, Vayu
  \item Triguna, Tridosha, Trivarga
\end{parts}

\medskip


\noindent   \textbf{Question 13.} What is the primary treatment for mental doshas in Ayurveda?\pts{1}

\begin{parts}
  \item Herbal remedies
  \item Meditation and self-awareness
  \item Detoxification therapies
  \item Dietary adjustments
\end{parts}

\medskip


\noindent   \textbf{Question 14.} Which type of Anubhava (experience) represents the real perception of object?\pts{1}

\begin{parts}
  \item Ayathartha Anubhava
  \item Yathartha Anubhava
  \item Pratyaksha Anubhava
  \item Yukti Anubhava
\end{parts}

\medskip


\noindent   \textbf{Question 15.} Which type of Agni is considered balanced and ideal?\pts{1}

\begin{parts}
  \item Tikshna Agni
  \item Vishama Agni
  \item Sama Agni
  \item Manda Agni
\end{parts}

\medskip


\noindent   \textbf{Question 16.} According to Ayurveda, which type of Agnibala (digestive capacity) is associated with Pitta Prakriti?\pts{1}

\begin{parts}
  \item Vishamagni
  \item Tikshnagni
  \item Mandagni
  \item Samagni
\end{parts}

\medskip


\noindent   \textbf{Question 17.} Which of the following is not included under Trayopastambha?\pts{1}

\begin{parts}
  \item Ahara
  \item Nidra
  \item Brahmacharya
  \item Vyayama
\end{parts}

\medskip


\noindent   \textbf{Question 18.} Which season is associated with accumulation (Sanchaya) of Kapha dosha?\pts{1}

\begin{parts}
  \item Shishira
  \item Vasanta
  \item Grishma
  \item Varsha
\end{parts}

\medskip


\noindent   \textbf{Question 19.} Which one is called as 'Ubhayendriya'?\pts{1}

\begin{parts}
  \item Sharir
  \item Manas
  \item Atma
  \item Chakshu
\end{parts}

\medskip


\noindent   \textbf{Question 20.} Which one is counted as a Nastika Darshan?\pts{1}

\begin{parts}
  \item Charvaka Darshan
  \item Samkhya Darshan
  \item Yoga Darshan
  \item Vaisheshika Darshan
\end{parts}

\medskip


\noindent   \textbf{Question 21.} In Ayurveda, which is not considered as Dhatu in the human body?\pts{1}

\begin{parts}
  \item Rasa
  \item Mamsa
  \item Mutra
  \item Meda
\end{parts}

\medskip


\noindent   \textbf{Question 22.} In Ayurveda, how many types of Dhatu (tissue) are in the human body?\pts{1}

\begin{parts}
  \item Five
  \item Six
  \item Seven
  \item Four
\end{parts}

\medskip


\noindent   \textbf{Question 23.} Number of Gyanendriya:\pts{1}

\begin{parts}
  \item 6
  \item 5
  \item 10
  \item 8
\end{parts}

\medskip


\noindent   \textbf{Question 24.} Which is not a classical branch of Ayurveda?\pts{1}

\begin{parts}
  \item Kaya Chikitsa
  \item Shalya Tantra
  \item Radiology
  \item Shalakya Tantra
\end{parts}

\medskip


\noindent   \textbf{Question 25.} Which one of the following is not abode of diseases?\pts{1}

\begin{parts}
  \item Sharir
  \item Manas
  \item Atma
  \item None of the above
\end{parts}

\medskip


\noindent   \textbf{Question 26.} Number of Padartha as per Ayurveda are:\pts{1}

\begin{parts}
  \item 6
  \item 5
  \item 8
  \item 7
\end{parts}

\medskip


\noindent   \textbf{Question 27.} Which Muhurta is considered best for wake up in the morning?\pts{1}

\begin{parts}
  \item Nisha
  \item Brahma
  \item Rajani
  \item None of the above
\end{parts}

\medskip


\noindent   \textbf{Question 28.} The number of Kriyakala are:\pts{1}

\begin{parts}
  \item 3
  \item 6
  \item 7
  \item 8
\end{parts}

\medskip


\noindent   \textbf{Question 29.} Which Padartha is considered as "Cause of Decrease"?\pts{1}

\begin{parts}
  \item Vishesha
  \item Samanya
  \item Guna
  \item Karma
\end{parts}

\medskip


\noindent   \textbf{Question 30.} Marmas are the points where injury can result in:\pts{1}

\begin{parts}
  \item Death
  \item Pain
  \item Disability
  \item All of the above
\end{parts}

\medskip


\noindent   \textbf{Question 31.} The Chikitsya Purusha has:\pts{1}

\begin{parts}
  \item 6 Dhatu
  \item 5 Dhatu
  \item 10 Dhatu
  \item 8 Dhatu
\end{parts}

\medskip


\noindent   \textbf{Question 32.} Which branch of Ayurveda deals with Eye, Ear, Nose, Throat disorders?\pts{1}

\begin{parts}
  \item Kaya Chikitsa
  \item Shalya
  \item Rasayan
  \item Shalakya
\end{parts}

\medskip


\noindent   \textbf{Question 33.} Therapeutic Purusha is also known as:\pts{1}

\begin{parts}
  \item Eka Dhatvatmaka
  \item Shaddhatvatmaka
  \item Chaturvinshati Dhatvatmaka
  \item All of the above
\end{parts}

\medskip


\noindent   \textbf{Question 34.} Dhatu which carries oxygen and life to the body is:\pts{1}

\begin{parts}
  \item Rasa
  \item Rakta
  \item Mamsa
  \item Shukra
\end{parts}

\medskip


\noindent   \textbf{Question 35.} Primary function of Asthi Dhatu is:\pts{1}

\begin{parts}
  \item Dharana
  \item Preenana
  \item Jeevana
  \item Lepana
\end{parts}

\medskip


\noindent   \textbf{Question 36.} Which Upadhatu is derived from Rasa Dhatu?\pts{1}

\begin{parts}
  \item Stanya and Artava
  \item Kandara
  \item Vasa
  \item Kesha
\end{parts}

\medskip


\noindent   \textbf{Question 37.} Asava & Arishta are:\pts{1}

\begin{parts}
  \item Medicated Oils
  \item Natural fermented liquid
  \item Herbal granules
  \item None of the above
\end{parts}

\medskip


\noindent   \textbf{Question 38.} Which department of Ayurveda deals with Medico-legal aspects?\pts{1}

\begin{parts}
  \item Agad Tantra & Vyavhar Ayurveda
  \item Samhita & Sanskrit
  \item Siddhant Darshan
  \item Bal Roga
\end{parts}

\medskip


\noindent   \textbf{Question 39.} The main objective of Medicinal Chemistry department is to support:\pts{1}

\begin{parts}
  \item Research in Ayurveda
  \item Drug development in Ayurveda
  \item Phyto Analysis in Ayurveda
  \item All of the above
\end{parts}

\medskip


\noindent   \textbf{Question 40.} Kriyakala are used to understand:\pts{1}

\begin{parts}
  \item Yoga
  \item Progression of disease
  \item Surgery
  \item Food intake
\end{parts}

\medskip


\noindent   \textbf{Question 41.} The ultimate origin of Ayurveda is from:\pts{1}

\begin{parts}
  \item Brahma
  \item Daksha Prajapati
  \item Ashwani Kumar
  \item Indra
\end{parts}

\medskip


\noindent   \textbf{Question 42.} Abode of diseases is/are:\pts{1}

\begin{parts}
  \item Sharir
  \item Manas
  \item Both (1) & (2)
  \item None of the above
\end{parts}

\medskip


\noindent   \textbf{Question 43.} Dhanwantari Jayanti is being celebrated as:\pts{1}

\begin{parts}
  \item National Doctors Day
  \item National Medicine Day
  \item National Herbs Day
  \item National Ayurveda Day
\end{parts}

\medskip


\noindent   \textbf{Question 44.} Which of the following compendium is not in the group of Brihattrayi?\pts{1}

\begin{parts}
  \item Charak Samhita
  \item Sushrut Samhita
  \item Vagbhatta Samhita
  \item Bhava Prakasha
\end{parts}

\medskip


\noindent   \textbf{Question 45.} How many types of Sweda are mentioned in Ayurveda?\pts{1}

\begin{parts}
  \item 2
  \item 13
  \item 4
  \item 8
\end{parts}

\medskip


\noindent   \textbf{Question 46.} Purva-Karma includes:\pts{1}

\begin{parts}
  \item Pre-Panchkarma Procedures
  \item Pre-Yoga Procedures
  \item Pre-operative Procedures
  \item Post-operative Procedures
\end{parts}

\medskip


\noindent   \textbf{Question 47.} How many times one should brush their tooth as per Ayurveda?\pts{1}

\begin{parts}
  \item 2 Times
  \item 5 Times
  \item 8 Times
  \item 10 Times
\end{parts}

\medskip


\noindent   \textbf{Question 48.} Which one is not Manasika Dosha?\pts{1}

\begin{parts}
  \item Satva
  \item Rajas
  \item Tamas
  \item None of the above
\end{parts}

\medskip


\noindent   \textbf{Question 49.} Which one of the following is benefits of Ahara?\pts{1}

\begin{parts}
  \item Good Complexion (Varna)
  \item Disease (Roga)
  \item Nourishment (Pushti)
  \item Both (1) & (3)
\end{parts}

\medskip


\noindent   \textbf{Question 50.} How many Upastambha are mentioned in Ayurveda?\pts{1}

\begin{parts}
  \item 2
  \item 1
  \item 6
  \item 3
\end{parts}

\medskip


\noindent   \textbf{Question 51.} What is the main goal of Rasayana therapy in Ayurveda?\pts{1}

\begin{parts}
  \item To reduce body heat
  \item To enhance digestion
  \item To promote overall health and longevity
  \item To increase appetite
\end{parts}

\medskip


\noindent   \textbf{Question 52.} According to the National Education Policy (NEP) 2020, why is Yoga integrated into the educational curriculum?\pts{1}

\begin{parts}
  \item To promote competitive exams
  \item To enhance mindfulness and cognitive performance
  \item To replace physical education
  \item To focus solely on physical fitness
\end{parts}

\medskip


\noindent   \textbf{Question 53.} Long-term alcohol abuse increases the risk of which type of cancer?\pts{1}

\begin{parts}
  \item Lung cancer
  \item Skin cancer
  \item Liver cancer
  \item Colon cancer
\end{parts}

\medskip


\noindent   \textbf{Question 54.} Reconstructive Surgery of damaged nose is termed as:\pts{1}

\begin{parts}
  \item Rhinoplasty
  \item Lobuloplasty
  \item Otoplasty
  \item Blepharoplasty
\end{parts}

\medskip


\noindent   \textbf{Question 55.} Ayurveda is considered as Upa-veda of which of the following Veda?\pts{1}

\begin{parts}
  \item Rig-veda
  \item Saama-veda
  \item Yajur-veda
  \item Atharva-veda
\end{parts}

\medskip


\noindent   \textbf{Question 56.} Who among the following are considered divine physicians of the Vedic period?\pts{1}

\begin{parts}
  \item Brahma
  \item Indra
  \item Ashwini Kumar
  \item Prajapati
\end{parts}

\medskip


\noindent   \textbf{Question 57.} The procedure in Ayurveda involving 'controlled application of heat' is known as:\pts{1}

\begin{parts}
  \item Agnikarma
  \item Jalaukavacharana
  \item Ksharakarma
  \item None of the above
\end{parts}

\medskip


\noindent   \textbf{Question 58.} Modified technique of Ksharasutra therapy developed at Banaras Hindu University is named as:\pts{1}

\begin{parts}
  \item LIFT
  \item IFTAK
  \item VAAFT
  \item None of the above
\end{parts}

\medskip


\noindent   \textbf{Question 59.} What is Shalakya Tantra in Ayurveda?\pts{1}

\begin{parts}
  \item Study of gastrointestinal disorders
  \item Study of diseases occurring above the clavicle
  \item Study of bone-related disorders
  \item Study of reproductive health
\end{parts}

\medskip


\noindent   \textbf{Question 60.} How many purification procedures are included in Shodhana Chikitsa?\pts{1}

\begin{parts}
  \item Four
  \item Five
  \item Six
  \item Seven
\end{parts}

\medskip


\noindent   \textbf{Question 61.} In Ayurveda, which daily practice is recommended for preventing diseases affecting the eyes?\pts{1}

\begin{parts}
  \item Netra Prakshalana
  \item Nasya
  \item Karnapuran
  \item Virechan
\end{parts}

\medskip


\noindent   \textbf{Question 62.} What is the term for dietary guidelines for a pregnant woman in Ayurveda?\pts{1}

\begin{parts}
  \item Garbhini Paricharya
  \item Sutika Paricharya
  \item Dinacharya
  \item Ritucharya
\end{parts}

\medskip


\noindent   \textbf{Question 63.} Direct acknowledgment using sense organs is called as:\pts{1}

\begin{parts}
  \item Anuman
  \item Pratyaksha
  \item Yukti
  \item Aaptopadesha
\end{parts}

\medskip


\noindent   \textbf{Question 64.} Identification of a new chemical entity as a potential therapeutic agent (From Hit to Lead) is known as:\pts{1}

\begin{parts}
  \item Drug discovery
  \item Drug Analysis
  \item Both of the above
  \item None of the above
\end{parts}

\medskip


\noindent   \textbf{Question 65.} The para-surgical procedures in Ayurveda are termed as:\pts{1}

\begin{parts}
  \item Shastra-karma
  \item Anushastra-karma
  \item Shodhan-karma
  \item Panch-karma
\end{parts}

\medskip


\noindent   \textbf{Question 66.} Which stage of Shodhana Karma involves Deepana-Pachana, Snehana and Swedana?\pts{1}

\begin{parts}
  \item Pradhana Karma
  \item Purva Karma
  \item Paschat Karma
  \item Upakrama Karma
\end{parts}

\medskip


\noindent   \textbf{Question 67.} In 'Ashtanga Ayurveda', which specialty deals with child diseases?\pts{1}

\begin{parts}
  \item Kayachikitsa
  \item Rasayana
  \item Kaumarbhritya
  \item Shalakya
\end{parts}

\medskip


\noindent   \textbf{Question 68.} What does CPR stand for?\pts{1}

\begin{parts}
  \item Cardiac Passive Resuscitation
  \item Cardio Pulmonary Resuscitation
  \item Chest Pulmonary Revival
  \item Cardio Pulmonary Rescue
\end{parts}

\medskip


\noindent   \textbf{Question 69.} Pratyaksha, Anuman and Aaptopadesh are called as:\pts{1}

\begin{parts}
  \item Pramaan
  \item Pariksha
  \item Both (1) & (2)
  \item None of the above
\end{parts}

\medskip


\noindent   \textbf{Question 70.} Which dimension is not included by WHO in the definition of "Health"?\pts{1}

\begin{parts}
  \item Physical
  \item Mental
  \item Social
  \item Ritual
\end{parts}

\medskip

\vfill
\rule{\linewidth}{0.4pt}

\begin{center}\small
\href{https://bhu-exam.vercel.app/}{Click here to visit our website}\\[4pt]
\href{https://me-aryan.vercel.app/}{Click here to visit our contributor}
\end{center}

\end{document}
